\documentclass{article}
\usepackage[T1]{fontenc}
\usepackage{lmodern}
\usepackage{graphicx}
\usepackage{amsmath,amssymb}
\usepackage[a4paper,margin=2cm]{geometry}
\usepackage[absolute,overlay]{textpos}
\setlength{\TPHorizModule}{1mm}
\setlength{\TPVertModule}{1mm}
\usepackage[indonesian]{babel}
\usepackage{hyperref}
\setlength{\emergencystretch}{2em}
% unit: Halaman 2

\begin{document}

% === Halaman 2 (llm) ===

\section*{Pendahuluan}

\begin{itemize}
    \item \textcolor{red}{Unbreakable cipher} merupakan klaim yang dibuat oleh kriptografer terhadap algoritma kriptografi yang dirancangnya.
    \item Namun, kebanyakan algoritma yang sudah pernah dibuat orang adalah \textcolor{red}{breakable cipher}.
    \item Caesar Cipher, Vigenere Cipher, Playfair Cipher, Affine Cipher, Enigma Cipher, Hill Cipher, dll sudah kadaluarsa karena \textcolor{red}{breakable cipher}.
\end{itemize}

\end{document}
